\documentclass[11pt,a4paper]{article}
% Preambulo comun para la memoria de Practica 2.
\usepackage[utf8]{inputenc}
\usepackage[T1]{fontenc}
\usepackage[spanish,es-tabla]{babel}
\usepackage[a4paper,margin=2.5cm]{geometry}

\usepackage{graphicx}
\usepackage{float}
\usepackage{caption}
\usepackage{subcaption}
\usepackage{booktabs}
\usepackage{array}
\usepackage{tabularx}
\usepackage{multirow}

\usepackage{amsmath}
\usepackage{amssymb}

\usepackage{xcolor}
\usepackage{listings}
\usepackage{hyperref}
\hypersetup{
  colorlinks=true,
  linkcolor=blue!50!black,
  urlcolor=blue!50!black,
  citecolor=blue!50!black,
}

\usepackage{enumitem}
\setlist{nosep}

\definecolor{codebg}{RGB}{245,245,245}
\definecolor{codekw}{RGB}{0,90,140}
\definecolor{codestr}{RGB}{160,40,40}
\definecolor{codecmt}{RGB}{90,120,90}
\lstdefinestyle{compact}{
  backgroundcolor=\color{codebg},
  basicstyle=\ttfamily\footnotesize,
  keywordstyle=\color{codekw}\bfseries,
  stringstyle=\color{codestr},
  commentstyle=\color{codecmt}\itshape,
  breaklines=true,
  showstringspaces=false,
  frame=single,
  framerule=0pt,
  xleftmargin=4pt,
  xrightmargin=4pt,
  aboveskip=4pt,
  belowskip=4pt,
}
\lstset{style=compact}

\newcommand{\fig}[3][0.85]{%
  \begin{figure}[H]
    \centering
    \includegraphics[width=#1\linewidth]{figures/#2}
    \caption{#3}
    \label{fig:#2}
  \end{figure}%
}


\title{Practica 2 -- Navegacion Autonoma con NAV2 en ROS\,2 Humble}
\author{
  [Nombre 1 Apellidos] -- [NIA 1] \\
  [Nombre 2 Apellidos] -- [NIA 2] \\
  [Nombre 3 Apellidos] -- [NIA 3] \\
  \\
  Master en Inteligencia Artificial -- Robotica Inteligente
}
\date{[Fecha de entrega]}

\begin{document}

\maketitle
\thispagestyle{empty}

\section{Introduccion y objetivo}
\label{sec:intro}

% TODO: 1 parrafo.
% - Construir el stack completo de navegacion en ROS 2 con NAV2: SLAM, EKF,
%   AMCL, costmaps, planner y controller, validado en Gazebo sobre el Husky.
% - Justificar uso de Husky A200 (Clearpath) por continuidad con el TFM
%   de inspeccion de torres eolicas.
% - Mencionar que la rama practica-2 contiene el repositorio completo del
%   TFM y que la memoria se centra en el scope de NAV2.
[Resumen del objetivo: integrar NAV2 sobre el Husky A200 simulado en Gazebo,
mapear el entorno con SLAM Toolbox, fusionar IMU + odometria con un EKF,
localizar globalmente con AMCL y completar misiones de navegacion enviando
metas desde RViz o ejecutando waypoints predefinidos. El proyecto en el que se
integra esta practica es la inspeccion autonoma de torres eolicas; en la rama
\texttt{practica-2} conviven otros paquetes del TFM (brazo, percepcion,
comportamiento) que no son objeto de evaluacion en esta entrega.]

\section{Arquitectura del sistema}
\label{sec:arquitectura}

% TODO: diagrama de bloques.
% - Gazebo -> /scan (via pointcloud_to_laserscan), /imu/data, /odom de ruedas
% - robot_localization (EKF) -> /odometry/filtered, TF odom->base_link
% - AMCL -> TF map->odom
% - nav2_bringup -> planner_server (NavFn), controller_server (RPP),
%   smoother_server, bt_navigator, behavior_server
% - lifecycle_manager orquesta el ciclo de vida
% - RViz para visualizar y enviar metas

\fig[0.9]{architecture_diagram.png}{Arquitectura NAV2 implementada sobre el Husky simulado.}

\section{Configuracion de launchers}
\label{sec:launchers}

% TODO: explicar los 4 launchers usados y por que existen separados.

\begin{table}[H]
  \centering
  \caption{Launchers utilizados en la Practica 2.}
  \label{tab:launchers}
  \begin{tabularx}{\linewidth}{lX}
    \toprule
    Launcher & Responsabilidad \\
    \midrule
    \texttt{simulation.launch.py}      & Gazebo Sim + Husky (Clearpath) + bridges + tuneo CycloneDDS \\
    \texttt{slam.launch.py}            & SLAM Toolbox (sync) + RViz con \texttt{slam.rviz} \\
    \texttt{navigation.launch.py}      & NAV2 sin AMCL, con EKF + SLAM online \\
    \texttt{navigation\_amcl.launch.py}& NAV2 + AMCL + mapa estatico ya generado \\
    \bottomrule
  \end{tabularx}
\end{table}

% TODO: subseccion sobre buenas practicas:
% - get_package_share_directory
% - DeclareLaunchArgument(use_sim_time, map, params_file)
% - lifecycle_manager con autostart=True
% - CYCLONEDDS_URI a tools/cyclonedds.xml (sube limite DDS de 64 a 240)
% - Relays /robot/tf -> /tf y /robot/robot_description -> /robot_description
%   por coexistencia de namespace Clearpath con NAV2.

\section{Mapeo con SLAM Toolbox}
\label{sec:slam}

% TODO:
% - slam_params.yaml en modo sync (sync_slam_toolbox_node).
% - Topicos consumidos: /scan, /tf, /odom.
% - Topicos publicados: /map, /map_metadata, /pose.
% - Conducir el robot manualmente con ps5_teleop o teleop_twist_keyboard.
% - Guardar el mapa con map_saver_cli -> ros2_ws/maps/wind_tower.{pgm,yaml,
%   posegraph,data}.
% Capturas:
% - figures/slam_in_progress.png  (RViz con mapa parcial)
% - figures/map_final.png         (mapa guardado)

\fig[0.8]{slam_in_progress.png}{Mapeo SLAM en progreso: ocupacion, escaneo LIDAR y trayectoria.}

\fig[0.55]{map_final.png}{Mapa final \texttt{wind\_tower.pgm} obtenido tras el recorrido.}

\section{Fusion sensorial con filtro de Kalman extendido}
\label{sec:ekf}

% TODO: explicar ekf_map.yaml:
% - Vector de estado: [x, y, theta, vx, omega].
% - two_d_mode: true (robot planar).
% - odom0 y imu0 con sus matrices _config:
%   * odom0 aporta pose (x, y, theta) y velocidades lineales.
%   * imu0 aporta orientacion absoluta y velocidad angular.
% - Frecuencia 30 Hz, sensor_timeout 0.1.
% - Salida: odom -> base_link y /odometry/filtered.
% - Mencionar el ajuste de covarianza Z (0.10 -> 0.03) para reflejar que el
%   robot es planar.

\begin{equation}
  \mathbf{x} = \begin{bmatrix} x & y & \theta & v_x & \omega \end{bmatrix}^T
\end{equation}

\begin{equation}
  \mathbf{x}_{k+1} = \mathbf{x}_k + \begin{bmatrix} v_k \cos\theta_k \\ v_k \sin\theta_k \\ \omega_k \\ 0 \\ 0 \end{bmatrix} \Delta t + \mathbf{w}_k
\end{equation}

\fig[0.6]{ekf_filtered_odom.png}{Comparativa entre \texttt{/odom} (sin fusion) y \texttt{/odometry/filtered} (EKF).}

\section{Localizacion con AMCL}
\label{sec:amcl}

% TODO: explicar amcl_params.yaml:
% - min_particles=500, max_particles=2000 (KLD-sampling).
% - alpha1..alpha4: ruido del modelo de movimiento.
% - laser_z_hit, laser_sigma_hit: modelo likelihood field.
% - update_min_d=0.25, update_min_a=0.2.

\begin{figure}[H]
  \centering
  \begin{subfigure}{0.48\linewidth}
    \centering
    \includegraphics[width=\linewidth]{figures/amcl_initial.png}
    \caption{Inicialmente: particulas dispersas.}
  \end{subfigure}
  \begin{subfigure}{0.48\linewidth}
    \centering
    \includegraphics[width=\linewidth]{figures/amcl_converged.png}
    \caption{Tras varios escaneos: convergencia.}
  \end{subfigure}
  \caption{Distribucion de particulas AMCL antes y despues de la convergencia.}
  \label{fig:amcl}
\end{figure}

\section{Mapas, costmaps y planner}
\label{sec:costmaps}

% TODO:
% - global_costmap: static_layer + obstacle_layer + inflation_layer.
% - local_costmap: rolling window, obstacle_layer + inflation_layer.
% - Planner: NavFn (Dijkstra).
% - Behavior tree por defecto.

\fig[0.85]{costmaps_rviz.png}{Costmap global (azul) y local (rojo) sobre el mapa estatico.}

\section{Tuning iterativo del controller}
\label{sec:tuning}

El controller local es \texttt{RegulatedPurePursuitController}. La rampa de 13°
presente en el mundo nos obligo a iterar la configuracion en dos fases, y este
proceso constituye una de las aportaciones mas relevantes de esta entrega.

\paragraph{Hipotesis 1 (descartada): \emph{"mas velocidad y aceleracion para vencer la pendiente"}.}
Subimos \texttt{desired\_linear\_vel} a 0.5\,m/s, \texttt{max\_accel} a 2.5 y
\texttt{max\_velocity} a 0.6. El razonamiento era que el Husky debia entrar con
momento suficiente para coronar la rampa. En la practica el robot patinaba en
los giros sobre la rampa, AMCL perdia coherencia con el LIDAR y la odometria
daba saltos.

\paragraph{Hipotesis 2 (actual): \emph{"aceleraciones suaves para no perder traccion"}.}
El problema dominante era traccion, no potencia. Bajamos las aceleraciones del
\texttt{velocity\_smoother} de 2.5 a 0.75 lineal y de 1.0 a 0.8 angular, y
suavizamos el seguimiento del path con un lookahead mas largo.

\begin{table}[H]
  \centering
  \caption{Parametros clave del controller antes y despues del tuning.}
  \label{tab:tuning}
  \begin{tabularx}{\linewidth}{lXrr}
    \toprule
    Parametro & Motivacion del cambio & Antes & Despues \\
    \midrule
    \texttt{lookahead\_dist}                  & Horizonte mas largo amortigua el seguimiento (anti-bandazos) & 0.9 & 1.2 \\
    \texttt{max\_lookahead\_dist}             & Mas margen en rectas largas       & 1.5 & 2.0 \\
    \texttt{min\_lookahead\_dist}             & Evita carrot point muy cerca       & 0.4 & 0.6 \\
    \texttt{rotate\_to\_heading\_angular\_vel}& Giros in-situ mas suaves           & 0.7 & 0.5 \\
    \texttt{rotate\_to\_heading\_min\_angle}  & Menos giros in-situ innecesarios   & 45° & 60° \\
    \texttt{max\_angular\_accel}              & Reduce jerk angular                & 2.0 & 1.2 \\
    \texttt{velocity\_smoother.max\_accel}    & Anti-slip en rampa (lineal)        & 2.5 & 0.75 \\
    \texttt{min\_approach\_linear\_velocity}  & Aproximacion suave sin derrape     & 0.15 & 0.06 \\
    \bottomrule
  \end{tabularx}
\end{table}

\paragraph{Lecciones aprendidas.}
\begin{enumerate}
  \item El problema en pendiente no era falta de potencia sino traccion: mas velocidad amplifica el slip.
  \item El controlador es muy sensible al horizonte de lookahead. Valores pequenos amplifican ruido del path; valores grandes lo amortiguan.
  \item La aceleracion del \texttt{velocity\_smoother} debe ser coherente con el resto de la cadena (smoother -> twist\_mux Clearpath -> diff-drive). Si una etapa satura, las demas no compensan.
  \item La hipotesis equivocada inicial nos costo tiempo pero genero un dataset claro de comparacion antes/despues que se ve en las metricas.
\end{enumerate}

\fig[0.85]{nav_full_run.png}{Captura combinada Gazebo + RViz tras el tuning: trayectoria suave y AMCL estable durante la subida.}

\section{Cadena TF}
\label{sec:tf}

\fig[0.75]{tf_frames.png}{Arbol TF obtenido con \texttt{ros2 run tf2\_tools view\_frames}.}

% TODO: explicar map -> odom (AMCL), odom -> base_link (EKF),
% base_link -> sensores (robot_state_publisher).

\section{Pruebas y resultados cuantitativos}
\label{sec:pruebas}

% TODO: rellenar la tabla de metricas tras las pruebas.

\begin{table}[H]
  \centering
  \caption{Metricas cuantitativas tras el tuning del controller.}
  \label{tab:metricas}
  \begin{tabularx}{\linewidth}{lXl}
    \toprule
    Metrica & Definicion & Valor \\
    \midrule
    Tiempo SLAM completo        & Recorrer todo el mundo                       & [pendiente] \\
    Convergencia AMCL           & Desde "2D Pose Estimate" hasta nube compacta & [pendiente] \\
    RMSE trayectoria            & Pose AMCL vs ground truth de Gazebo          & [pendiente] \\
    Tasa replanificacion        & Eventos/min con obstaculo dinamico           & [pendiente] \\
    Aborts del BT en rampa      & Antes y despues del tuning                   & [pendiente] \\
    \bottomrule
  \end{tabularx}
\end{table}

\section{Problemas encontrados y soluciones}
\label{sec:problemas}

% TODO: vinetas con problemas reales. Ejemplos plausibles:
% - DDS slots agotados (Clearpath + NAV2 + bridges) -> CYCLONEDDS_URI con
%   file://tools/cyclonedds.xml (sube limite de 64 a 240).
% - LIDAR 3D Velodyne -> requiere pointcloud_to_laserscan para alimentar
%   SLAM/AMCL en 2D.
% - QoS incompatible en /scan -> scan_qos_bridge.
% - Spawn y orientacion del robot mirando al sur (yaw=3.14) para entrar al tubo.
% - Slip en rampa -> ver tuning seccion.

\begin{itemize}
  \item [Problema 1]: [descripcion y solucion].
  \item [Problema 2]: [descripcion y solucion].
  \item [Problema 3]: [descripcion y solucion].
\end{itemize}

\section{Conclusiones}
\label{sec:conclusiones}

% TODO: 1 parrafo final con lo aprendido + integracion con el TFM
% (Practicas 3 y 4 reutilizan este stack de navegacion).
[Conclusiones breves: NAV2 funciona end-to-end sobre el Husky simulado, listo
para integrarse con el resto del TFM. El tuning iterativo del controller deja
clara la diferencia entre el modelo mental inicial (mas potencia) y el
verdadero limite del sistema (traccion).]

\end{document}
